% ============================================================
\section{Analysis of Results}
\label{sec:discussion-analysis}
% ============================================================

The experimental results validate the core hypothesis of this thesis: that a
\textbf{dual-signal anomaly detection framework} combining a statistical
Z-score detector with a machine learning IsolationForest model produces more
reliable incident signals than either detector in isolation. The key evidence
is the 94.9\% model agreement rate combined with a 21.1\% mean Jaccard
similarity. This pattern is consistent with the theoretical complementarity of
the two signals: Z-score operates on a single dimension (\texttt{heavy\_ratio}
deviation from a per-route, per-time-slot baseline), whereas IsolationForest
operates on a seven-dimensional feature space that includes multivariate
congestion ratios and cyclical time components. An event must therefore exhibit
\textit{both} a statistical rarity and a multivariate distributional
abnormality to trigger a joint alert, significantly reducing false positive
exposure.

The Medallion Lakehouse architecture played a critical supporting role. The
separation of Bronze (raw), Silver (cleaned), and Gold (aggregated) layers
allowed the batch pipeline to evolve the schema of downstream tables
independently of the raw ingestion stream. In practice, the \texttt{source}
column was added to the Silver schema post-deployment without any reprocessing
of Bronze data, demonstrating the schema-evolution capability of Apache
Iceberg~\cite{iceberg_paper}.

The Welford online algorithm~\cite{welford1962} proved essential for the
real-time Z-score computation. Its $O(1)$ memory footprint ensured that the
Online Service remained stable over the 40-day observation period without
memory growth, even at 5-minute polling intervals across 20 concurrent routes.
The crash-recovery mechanism --- reconstructing the Welford accumulator from
stored PostgreSQL statistics --- was validated during two service restarts
without observable discontinuities in the per-route rolling statistics.

% ============================================================
\section{Strengths}
\label{sec:strengths}
% ============================================================

\textbf{Streaming-first, low-latency design.}
The p95 end-to-end ingestion lag below 20~seconds is a significant improvement
over batch-based systems that typically report latencies in the range of
5--30~minutes~\cite{batch_latency_survey}. The Server-Sent Events delivery
model eliminates polling overhead and provides a push-based update channel
that scales linearly with the number of frontend clients.

\textbf{Unified data platform.}
By combining a streaming path (Redpanda $\to$ PostgreSQL) and a batch path
(MinIO $\to$ Iceberg Silver/Gold) within the same Docker Compose stack,
UrbanPulse avoids the complexity of a Lambda architecture while retaining
the analytical depth of a Lakehouse. The unified Nessie catalog enables
time-travel queries over Iceberg tables, facilitating retrospective
investigation of past anomaly windows.

\textbf{Explainable anomaly detection.}
The RAG-enriched LLM pipeline provides natural-language explanations that
contextualise anomaly scores with route geography, historical patterns, and
live weather data. This capability directly addresses the \textit{semantic
gap} identified in the problem statement: raw anomaly scores are opaque to
operators, whereas a two-sentence explanation referencing specific road names
and congestion ratios is immediately actionable.

\textbf{Operational robustness.}
The system has operated continuously in production at
\texttt{tyr1on.io.vn} with automated CI/CD (ruff lint, mypy strict, unit
tests on every commit) and Grafana/Loki observability. The 30-minute
Telegram alert cooldown and the \texttt{both\_anomaly} conjunction rule
prevent alert fatigue, a common failure mode in production monitoring systems.

% ============================================================
\section{Limitations}
\label{sec:limitations}
% ============================================================

\textbf{Absence of ground-truth labels.}
The most significant limitation is the lack of labelled anomaly events. The
proxy metrics (model agreement, Jaccard similarity) measure internal
consistency between the two detectors rather than true recall against real
traffic incidents. A dataset of confirmed incidents from the Ho Chi Minh City
Department of Transport~\cite{hcmc_transport_report} would enable computation
of true precision and recall.

\textbf{VietMap API scope restriction.}
Traffic data is sourced exclusively from the VietMap Directions API, which
provides route-level congestion ratios but does not expose individual vehicle
counts, speed distributions, or incident reports. The congestion signal is
therefore an indirect proxy for underlying traffic density. Events such as
flooding, road works, or accidents are not distinguishable from general
high-demand congestion without supplementary data sources.

\textbf{IsolationForest model staleness.}
Models are retrained every 6~hours on historical Gold data. If an anomalous
pattern persists for more than 6~hours (e.g., a multi-day road closure), the
model may adapt and cease flagging the event as anomalous. This
\textit{concept drift} problem~\cite{chandola2009anomaly} is partially
mitigated by the Z-score signal, which continues to compare against a more
slowly updated baseline.

\textbf{Cold-start and low-variance routes.}
Four suburban routes (zone4 outbound, zone4-to-zone2, zone5 outbound) showed
zero Z-score anomaly events over the 40-day collection period. In low-variance
traffic conditions, the baseline standard deviation approaches zero, making
the Z-score numerically unstable. The implementation applies a minimum stddev
floor of \texttt{mean * 0.1} in these cases, but this is a heuristic that
may generate spurious signals immediately after the baseline is first computed.

\textbf{Hardware constraints.}
The platform is deployed on a single-node homelab (Intel i5, 8~GB VRAM for
GPU inference). Horizontal scaling of the Serving API and Batch Service to
multiple nodes was not evaluated. The Redpanda broker operates as a
single-partition configuration; a multi-partition, multi-replica setup
would be required for fault-tolerant production deployments at city scale.

% ============================================================
\section{Comparison with Related Approaches}
\label{sec:comparison}
% ============================================================

Table~\ref{tab:comparison} compares UrbanPulse against three architectural
patterns discussed in the Related Work chapter.

\begin{table}[htbp]
\centering
\caption{Comparison of UrbanPulse against reference architectures.}
\label{tab:comparison}
\resizebox{\textwidth}{!}{%
\begin{tabular}{lcccc}
\hline
\textbf{Feature} &
\textbf{UrbanPulse} &
\textbf{Batch-only} &
\textbf{Lambda Arch.} &
\textbf{Kappa Arch.} \\
\hline
Ingestion latency (p95)   & $<$20 s & 5--30 min & $<$1 min & $<$1 min \\
Schema evolution           & \cmark (Iceberg) & \xmark & \xmark & Partial \\
Time-travel queries        & \cmark (Nessie)  & \cmark & \xmark & \xmark \\
Dual-signal detection      & \cmark & \xmark & \xmark & \xmark \\
Explainability (LLM/RAG)   & \cmark & \xmark & \xmark & \xmark \\
Single codebase            & \cmark & \cmark & \xmark & \cmark \\
Ground-truth evaluation    & \xmark & \cmark & Partial & Partial \\
\hline
\end{tabular}%
}
\end{table}

\noindent UrbanPulse competes favourably on latency and feature richness
against traditional batch-only and Lambda architectures. Its primary
disadvantage is the absence of ground-truth evaluation, which is shared with
most deployed streaming anomaly systems that operate on unlabelled
real-world data.

% ============================================================
\section{Future Work}
\label{sec:future-work}
% ============================================================

Several extensions would strengthen the platform and address its current
limitations:

\begin{itemize}
  \item \textbf{Ground-truth label collection}: Integrate with the HCMC
        Department of Transport incident API to build a labelled dataset and
        compute precision, recall, and F1 for both detectors.
  \item \textbf{Multi-source ingestion}: Add vehicle loop detector data and
        social media signal (Twitter/X flood reports) as additional Kafka
        topics to improve incident classification.
  \item \textbf{Online learning for IsolationForest}: Replace the 6-hour
        batch retrain with an incremental Half-Space Trees (HST)
        detector~\cite{chandola2009anomaly} to reduce concept drift exposure.
  \item \textbf{Multi-node deployment}: Evaluate Kubernetes-based deployment
        with horizontal pod autoscaling for the Serving API to support
        thousands of concurrent SSE clients.
  \item \textbf{Federated RAG}: Extend the ChromaDB collections to include
        meteorological event reports, public holiday calendars, and
        construction permit data for richer root-cause analysis.
\end{itemize}
